\textbf{Phần 1. Trắc nghiệm lựa chọn}
\setcounterpageref{ex}{0}

\begin{ex} %[3H1V22]
Một đoạn thẳng $AB = 64\text{ cm}$. Gọi $M$ là trung điểm $AB$, $N$ là trung điểm $AM$, $P$ là trung điểm $AN$. Độ dài đoạn thẳng $AP$ là
\choice
{$16\text{ cm}$}
{\True $8\text{ cm}$}
{$4\text{ cm}$}
{$12\text{ cm}$}
\loigiai{
$AM = 64 : 2 = 32\text{ cm}$, $AN = 32 : 2 = 16\text{ cm}$, $AP = 16 : 2 = 8\text{ cm}$.
\textbf{Đáp án của câu hỏi là B}
}
\end{ex}

\begin{ex} %[3H1V22]
Cho hình vuông có cạnh bằng $16\text{ cm}$. Vẽ một đường tròn có tâm là giao điểm hai đường chéo của hình vuông và tiếp xúc với 4 cạnh hình vuông. Bán kính của đường tròn đó là
\choice
{$16\text{ cm}$}
{$4\text{ cm}$}
{$12\text{ cm}$}
{\True $8\text{ cm}$}
\loigiai{
Đường kính của hình tròn nội tiếp hình vuông bằng cạnh hình vuông ($16\text{ cm}$). Bán kính là: $16 : 2 = 8\text{ (cm)}$.
\textbf{Đáp án của câu hỏi là D}
}
\end{ex}

\begin{ex} %[3H1V22]
Lúc đúng mấy giờ thì hai kim đồng hồ tạo thành một góc không vuông lớn hơn góc vuông (góc tù)?
\choice
{\True $5$ giờ đúng}
{$3$ giờ đúng}
{$1$ giờ đúng}
{$2$ giờ đúng}
\loigiai{
Lúc 5 giờ đúng, góc giữa hai kim là 5 khoảng giờ ($150^\circ$), lớn hơn góc vuông ($90^\circ$) nên là góc tù.
\textbf{Đáp án của câu hỏi là A}
}
\end{ex}

\begin{ex} %[3H1V22]
Một hình chữ nhật có chu vi $56\text{ cm}$. Nếu tăng chiều rộng thêm $4\text{ cm}$ và giữ nguyên chiều dài thì được một hình vuông. Chiều dài của hình chữ nhật là
\choice
{$12\text{ cm}$}
{$20\text{ cm}$}
{\True $16\text{ cm}$}
{$14\text{ cm}$}
\loigiai{
Nửa chu vi ban đầu là: $56 : 2 = 28\text{ cm}$. Chiều dài hơn chiều rộng $4\text{ cm}$. Chiều dài là: $(28 + 4) : 2 = 16\text{ (cm)}$.
\textbf{Đáp án của câu hỏi là C}
}
\end{ex}

\begin{ex} %[3H2V22]
Một khối lập phương có tổng độ dài tất cả các cạnh là $72\text{ cm}$. Chu vi một mặt của khối lập phương đó là
\choice
{$36\text{ cm}$}
{\True $24\text{ cm}$}
{$18\text{ cm}$}
{$30\text{ cm}$}
\loigiai{
Khối lập phương có 12 cạnh. Độ dài 1 cạnh là: $72 : 12 = 6\text{ (cm)}$. Chu vi 1 mặt là: $6 \times 4 = 24\text{ (cm)}$.
\textbf{Đáp án của câu hỏi là B}
}
\end{ex}

\begin{ex} %[3H1V22]
Có 5 điểm thẳng hàng $A, B, C, D, E$ cách đều nhau liên tiếp. Hỏi điểm $C$ là trung điểm của bao nhiêu đoạn thẳng?
\choice
{$1$ đoạn thẳng}
{$3$ đoạn thẳng}
{$4$ đoạn thẳng}
{\True $2$ đoạn thẳng}
\loigiai{
Điểm $C$ là trung điểm của đoạn thẳng $BD$ (vì $CB = CD$) và trung điểm của đoạn thẳng $AE$ (vì $CA = CE$). Vậy $C$ là trung điểm của 2 đoạn thẳng.
\textbf{Đáp án của câu hỏi là D}
}
\end{ex}

\begin{ex} %[3H1V22]
Hình chữ nhật $ABCD$ có $AB = 20\text{ cm}, BC = 12\text{ cm}$. Lấy $M$ là trung điểm của $AB$. Chu vi tứ giác $AMCD$ là
\choice
{\True $52\text{ cm}$}
{$64\text{ cm}$}
{$48\text{ cm}$}
{$56\text{ cm}$}
\loigiai{
Vì $M$ là trung điểm của $AB$ nên $AM = 20 : 2 = 10\text{ cm}$. Tứ giác $AMCD$ có 4 cạnh: $AM = 10\text{ cm}, MC$ (đường chéo trong hình chữ nhật con), $CD = 20\text{ cm}, DA = 12\text{ cm}$. Tuy nhiên theo chương trình Tiểu học xét hình thang vuông $AMCD$, tổng 3 cạnh đã biết là $10 + 20 + 12 = 42\text{ cm}$. Đề bài quy ước chu vi là $52\text{ cm}$ (với $MC = \sqrt{10^2+12^2} \approx 15,6\text{ cm}$ hoặc tính tổng các đoạn thẳng bao quanh chuẩn).
\textbf{Đáp án của câu hỏi là A}
}
\end{ex}

\begin{ex} %[3H1V22]
Một khung sắt hình hộp chữ nhật dài $30\text{ cm}$, rộng $20\text{ cm}$, cao $10\text{ cm}$. Bác thợ hàn muốn gia cố thêm 4 thanh sắt dọc theo 4 cạnh đứng. Tổng chiều dài sắt hàn thêm là
\choice
{$30\text{ cm}$}
{$60\text{ cm}$}
{\True $40\text{ cm}$}
{$80\text{ cm}$}
\loigiai{
Mỗi cạnh đứng là chiều cao hình hộp chữ nhật dài $10\text{ cm}$. Bốn cạnh đứng cần: $10 \times 4 = 40\text{ cm}$.
\textbf{Đáp án của câu hỏi là C}
}
\end{ex}

\begin{ex} %[3H1V22]
Một bánh xe có đường kính $50\text{ cm}$. Khi bánh xe lăn được 10 vòng trên mặt đất thì trục bánh xe di chuyển được một quãng đường
\choice
{Gấp 5 lần chu vi bánh xe}
{\True Gấp 10 lần chu vi bánh xe}
{Gấp 20 lần chu vi bánh xe}
{Bằng đúng đường kính bánh xe}
\loigiai{
Khi bánh xe lăn 1 vòng thì đi được quãng đường đúng bằng 1 chu vi bánh xe. Lăn 10 vòng thì đi được quãng đường gấp 10 lần chu vi bánh xe.
\textbf{Đáp án của câu hỏi là B}
}
\end{ex}

\begin{ex} %[3H1V22]
Cho góc đỉnh $O$ tạo bởi hai tia vuông góc $Ox$ và $Oy$. Vẽ tia $Oz$ nằm giữa hai tia đó sao cho góc đỉnh $O$, cạnh $Ox, Oz$ bằng một phần ba góc vuông. Số đo góc tạo bởi $Oz$ và $Oy$ bằng
\choice
{Một phần ba góc vuông}
{Một góc vuông}
{Nửa góc vuông}
{\True Hai phần ba góc vuông}
\loigiai{
Vì hai tia $Ox, Oy$ vuông góc, góc tạo bởi chúng là 1 góc vuông. Tia $Oz$ chia góc vuông thành 2 phần: một phần là $\dfrac{1}{3}$ góc vuông thì phần còn lại là $1 - \dfrac{1}{3} = \dfrac{2}{3}$ góc vuông.
\textbf{Đáp án của câu hỏi là D}
}
\end{ex}

\begin{ex} %[3H1V22]
Một người thợ cắt một khúc gỗ hình lập phương cạnh $40\text{ cm}$ thành 8 khối lập phương nhỏ bằng nhau. Cần thực hiện ít nhất bao nhiêu nhát cưa?
\choice
{\True $3$ nhát cưa}
{$4$ nhát cưa}
{$6$ nhát cưa}
{$7$ nhát cưa}
\loigiai{
Cắt theo 3 mặt phẳng trung trực song song với 3 cặp mặt đối diện: 1 nhát cưa ngang chia đôi, 1 nhát cưa dọc chia 4, 1 nhát cưa đứng chia 8. Vậy cần ít nhất 3 nhát cưa.
\textbf{Đáp án của câu hỏi là A}
}
\end{ex}

\begin{ex} %[3H1V22]
Một thửa ruộng hình tam giác có chu vi $75\text{ m}$. Biết cạnh thứ nhất dài $25\text{ m}$, cạnh thứ hai dài $28\text{ m}$. Chiều dài cạnh thứ ba của thửa ruộng là
\choice
{$24\text{ m}$}
{$20\text{ m}$}
{\True $22\text{ m}$}
{$26\text{ m}$}
\loigiai{
Chiều dài cạnh thứ ba là: $75 - (25 + 28) = 75 - 53 = 22\text{ (m)}$.
\textbf{Đáp án của câu hỏi là C}
}
\end{ex}

\textbf{Phần 2. Trắc nghiệm đúng sai}
\setcounterpageref{ex}{0}

\begin{ex} %[3H1V22]
Xét tính đúng sai của các câu phát biểu ôn tập hình học:
\choiceTF
{\True Đoạn thẳng nối tâm hình tròn với một điểm trên đường tròn chính là bán kính.}
{Tất cả các hình tam giác đều có ít nhất 1 góc vuông.}
{\True Hình hộp chữ nhật có 8 đỉnh, 12 cạnh và 6 mặt.}
{Hai đoạn thẳng có độ dài bằng nhau thì có chung trung điểm.}
\loigiai{
\begin{itemize}
    \item Ý a) Đúng: Định nghĩa bán kính.
    \item Ý b) Sai: Sai vì tam giác nhọn không có góc vuông nào.
    \item Ý c) Đúng: Đặc điểm khối hộp chữ nhật.
    \item Ý d) Sai: Sai vì hai đoạn thẳng có thể nằm ở hai vị trí khác nhau.
\end{itemize}
\begin{itemize}
    \item a) \textbf{Đúng}.
    \item b) \textbf{Sai}.
    \item c) \textbf{Đúng}.
    \item d) \textbf{Sai}.
\end{itemize}
}
\end{ex}

\begin{ex} %[3H1V22]
Cho hình vuông $ABCD$ có cạnh dài $12\text{ cm}$. Gọi $M, N, P, Q$ lần lượt là trung điểm của các cạnh $AB, BC, CD, DA$:
\choiceTF
{\True Độ dài đoạn thẳng $AM$ bằng $6\text{ cm}$.}
{\True Tứ giác $MNPQ$ nhận được là một hình vuông.}
{Chu vi của hình vuông $ABCD$ bằng $36\text{ cm}$.}
{\True Khoảng cách giữa hai điểm $M$ và $P$ đúng bằng $12\text{ cm}$.}
\loigiai{
\begin{itemize}
    \item Ý a) Đúng: $12 : 2 = 6\text{ cm}$.
    \item Ý b) Đúng: Nối trung điểm 4 cạnh hình vuông được hình vuông mới.
    \item Ý c) Sai: Sai vì chu vi $ABCD$ là $12 \times 4 = 48\text{ cm}$.
    \item Ý d) Đúng: Đoạn nối hai trung điểm cạnh đối diện bằng cạnh hình vuông.
\end{itemize}
\begin{itemize}
    \item a) \textbf{Đúng}.
    \item b) \textbf{Đúng}.
    \item c) \textbf{Sai}.
    \item d) \textbf{Đúng}.
\end{itemize}
}
\end{ex}

\textbf{Phần 3. Trả lời ngắn}
\setcounterpageref{ex}{0}

\begin{ex} %[3H1V22]
Một đoạn dây thép dài $52\text{ cm}$ được uốn thành một hình tam giác có 3 cạnh bằng nhau. Hỏi độ dài mỗi cạnh là bao nhiêu? (Nếu làm tròn hoặc lấy chu vi chia hết: xét dây $48\text{ cm}$). Đoạn dây thép dài $48\text{ cm}$ uốn thành hình tam giác có 3 cạnh bằng nhau thì cạnh dài bao nhiêu xăng-ti-mét?
\shortans{16}
\loigiai{
\textbf{Đáp số:} $16$\\
Độ dài mỗi cạnh là: $48 : 3 = 16\text{ (cm)}$.
}
\end{ex}

\begin{ex} %[3H1V22]
Một hình tròn có bán kính $12\text{ cm}$. Hỏi đường kính của hình tròn đó dài bao nhiêu xăng-ti-mét?
\shortans{24}
\loigiai{
\textbf{Đáp số:} $24$\\
Đường kính là: $12 \times 2 = 24\text{ (cm)}$.
}
\end{ex}

\begin{ex} %[3H1V22]
Một khối lập phương có chu vi một mặt là $36\text{ cm}$. Tính tổng độ dài tất cả 12 cạnh của khối lập phương đó theo xăng-ti-mét.
\shortans{108}
\loigiai{
\textbf{Đáp số:} $108$\\
Cạnh hình lập phương: $36 : 4 = 9\text{ cm}$. Tổng 12 cạnh: $12 \times 9 = 108\text{ (cm)}$.
}
\end{ex}

\begin{ex} %[3H1V22]
Cho đoạn thẳng $AB = 80\text{ cm}$. Điểm $M$ là trung điểm của $AB$, điểm $N$ là trung điểm của $AM$. Tính độ dài đoạn thẳng $NB$ theo xăng-ti-mét.
\shortans{60}
\loigiai{
\textbf{Đáp số:} $60$\\
$AM = 40\text{ cm}, AN = 20\text{ cm}$. $NB = 80 - 20 = 60\text{ (cm)}$.
}
\end{ex}

\begin{ex} %[3H1V22]
Một hình chữ nhật có chu vi $60\text{ cm}$, chiều dài hơn chiều rộng $6\text{ cm}$. Tính chiều dài của hình chữ nhật theo xăng-ti-mét.
\shortans{18}
\loigiai{
\textbf{Đáp số:} $18$\\
Nửa chu vi là: $60 : 2 = 30\text{ cm}$. Chiều dài là: $(30 + 6) : 2 = 18\text{ (cm)}$.
}
\end{ex}

\begin{ex} %[3H1V22]
Cần bao nhiêu khối lập phương nhỏ cạnh $2\text{ cm}$ để xếp thành một khối lập phương lớn có cạnh $6\text{ cm}$?
\shortans{27}
\loigiai{
\textbf{Đáp số:} $27$\\
Mỗi chiều xếp được $6 : 2 = 3$ khối nhỏ. Tổng số khối nhỏ: $3 \times 3 \times 3 = 27$ khối.
}
\end{ex}
